%%%%%%%%%%%%%%%%%%%%%%%%%%%%%%%%%%%%%%%%%%%
\section{Tame fundamental groups of schemes}
\label{sec:pit schemes}
%%%%%%%%%%%%%%%%%%%%%%%%%%%%%%%%%%%%%%%%%%%

The goal of this section is to study the tame fundamental group of a connected scheme~$X$ relative to a base scheme $S$.
It is defined as the automorphism group of a fibre functor of the Galois category of tame covers of~$X$.
Accordingly, we start with a discussion of tame extensions of valued fields (\cref{sec:tame extensions valued fields}) and then of tame covers of schemes including a comparison of different notions of tameness found in the literature (\cref{sec:adic tame discrete case}).
The definition of tameness we will be using is the maximal one, which imposes tameness conditions on all valuations with a centre on $S$ of every residue field of~$X$. 
After a review of Galois categories in \cref{sec:Gal cats}, we are in the position to define the tame fundamental group in \cref{sec:tame pi1}.
The final results of this section are \cref{thm:naked pi1tame fg} and \cref{thm:naked pi1tame fp} showing that the tame fundamental group of a connected scheme of finite type over an algebraically closed field~$k$ is finitely generated and even finitely presented if it satisfies additional regularity assumptions.
We deduce these theorems from existing results \cite{EsnaultShustermanSrinivas2022:FinitePresentationTame} on finite generation and presentation of the tame fundamental group of a smooth variety with a projective snc compactification.
The reduction to the smooth case is done by a descent argument, for which we lay foundations in \cref{sec:descent for Gal cats}.

%%%---------------------------------------
\subsection{Tame extensions of valued fields} 
\label{sec:tame extensions valued fields}
%%%---------------------------------------

In this opening subsection we review tame extensions of valued fields, which form the basis of any discussion of tame phenomena in arithmetic geometry.
We give various characterisations of tameness and study the Galois group of the maximal tame extension.
Moreover, we explain how to describe tame extensions of compositions of valuations.

Throughout we will use the notation $(K,K^+)$ for a valued field with valuation ring~$K^+$ and write $K^\succ$ for the residue field of~$K^+$, which we also call the \textbf{specialisation field} of~$K$.
We denote the corresponding value group by~$\Gamma_K$, for which we use multiplicative notation, so the valuation of~$K$ is a multiplicative map
\[
 \lvert \cdot \rvert \colon K \longrightarrow \Gamma_K \cup \{0\}
\]
sending~$1$ to~$1$ and~$0$ to~$0$ and satisfying the strong triangle inequality.

For the valued field $(K,K^+)$ we fix an algebraic closure $(\overline{K},\overline{K}^+)$.
Then we have intermediate extensions of valued fields
\[
 (K,K^+) \subseteq (K^\rh,K^{\rh,+}) \subseteq (K^\sh,K^{\sh,+}) \subseteq (\overline{K},\overline{K}^+),
\]
where $(K^\rh,K^{\rh,+})$ is the {\bf henselisation}, and $(K^\sh,K^{\sh,+})$ is the {\bf strict henselisation} (or {\bf maximal unramified extension}).
Accordingly, $(K,K^+)$ is {\bf henselian} (resp.\ {\bf strictly henselian}) if and  only if $(K,K^+) = (K^\rh,K^{\rh,+})$ (resp.\ $(K,K^+) = (K^\sh,K^{\sh,+})$).

\begin{defi}
\label{defi:tame-extension-of-valued-fields}
    A finite extension of valued fields $(K,K^+) \to (L,L^+)$ is called \textbf{tame} if the residue characteristic does not divide the degree $[L^\sh:K^\sh]$.
    An algebraic extension of valued fields is \textbf{tame} if all of its finite subextensions are tame.
\end{defi}

From the definition it is clear that compositions and base changes of tame extensions are tame.
Consequently, there is a \textbf{maximal tamely ramified} extension $(K^t,K^{t,+})$ of $(K,K^+)$ (also called \textbf{tame henselisation}) complementing the above tower:
\[
 (K,K) \subseteq (K^\rh,K^{\rh,+}) \subseteq (K^\sh,K^{\sh,+}) \subseteq (K^t,K^{t,+}) \subseteq (\overline{K},\overline{K}^+).
\]
The valued field $(K,K^+)$ is called {\bf tamely henselian} if $(K,K^+) = (K^t,K^{t,+})$.
In the literature, the term {\bf tame closed} is often used for tamely henselian.

The structure of tame extensions is best understood via Galois theory.
Suppose that $(L,L^+)/(K,K^+)$ is a Galois extension of valued fields.
Its inertia group~$I$ is naturally identified with $\Gal(L^\sh/K^\sh)$.
We can now understand tame extensions in terms of the Kummer pairing
\[
 I \times (\Gamma_L/\Gamma_K) \longrightarrow \mu(L^\succ) \subseteq \bQ/\bZ'(1) 
\]
that maps $(\sigma,\gamma) \in I \times \Gamma_L$ to the residue class of $\sigma(a)/a$ in $(L^\succ)^\times$, where $a \in L$ is an element with $\lvert a \rvert = \gamma$.
It induces a surjective group homomorphism (the tame character)
\begin{equation} \label{Kummer-hom}
 I \longrightarrow \Hom(\Gamma_L/\Gamma_K,\bQ/\bZ'(1)), 
\end{equation}
whose kernel~$I_p$, called wild inertia, is a $p$-group, where $p \ge 0$ is the residue characteristic of~$K$ (see \cite[Proposition~6.2.12]{GabberRamero}).
Applying these considerations to the extension $K^t/K^\rh$, we obtain the following structural result for the Galois group of the maximal tame extension


\begin{lem} \label{lem:max tame extension}
    Suppose that $(K,K^+)$ is henselian of residue characteristic $p \ge 0$.
    \begin{enumerate}[ref=(\arabic*)]
        \item \label{lemitem:tame galois sequence}
        There is a short exact sequence
        \[
        1 \la \Hom(\Gamma_K,\hZ'(1)) \la \Gal(K^t/K) \la \Gal(K^\sh/K) \la 1
        \]
        where the map on the left comes from the Kummer theory pairings for $p \nmid n$:
        \[
        \Gal(K^t/K^\sh) \times \Gamma_{K^t}/\Gamma_K \la \bQ/\bZ'(1), 
        \qquad (\sigma,\gamma) \mapsto [\sigma(a)/a]
        \]
        with $a \in K^t$ of valuation $\gamma$, and $[\sigma(a)/a]$ is the residue class of $\sigma(a)/a$ in $(K^{t,\succ})^\times$.
  
        \item \label{lemitem:tame galois split}
        The sequence in \labelcref{lemitem:tame galois sequence} is split and 
        $\Gal(K^\sh/K)$ is canonically isomorphic to the absolute Galois group $\Gal_{K^\succ}$ of the residue field.
    \end{enumerate}
\end{lem}
\begin{proof}
    We first note that 
    \[
        \Hom(\Gamma_K,\hZ'(1)) = \varinjlim_{p \mid n} \Hom\big(\Gamma_K/n\Gamma_K,\mu_n(K^{t,\succ})\big) = \Hom(\Gamma_{K^t}/\Gamma_K,\bQ/\bZ'(1)).
    \]
     Since~$K$ is henselian, the extension $K^\sh/K$ is Galois with group isomorphic to $\Gal_{K^\succ}$.
    The kernel of the surjection $\Gal(K^t/K) \twoheadrightarrow \Gal(K^\sh/K) = \Gal_{K^\succ}$ is the inertia group of the extension $K^t/K$ by definition.
    Via the Kummer pairing it identifies with $\Hom(\Gamma_K,\hZ'(1))$ because the wild inertia group is trivial (see the discussion above).
    This accounts for the exact sequence in \labelcref{lemitem:tame galois sequence}.
    That the sequence splits follows from  \cite{MelnikovTavgen85} or  \cite[Theorem~2.2]{KuhlmannPankRoquette}.
\end{proof}  

\begin{cor}
\label{cor:pitame injective for fields}
    Let $K'/K$ be a finite extension of henselian valued fields. Then the induced map $\Gal(K'^t/K') \to \Gal(K^t/K)$ is an isomorphism onto an open subgroup.
\end{cor}

\begin{proof}
We compare both tame fundamental groups by means of the short exact sequence in \cref{lem:max tame extension}~\labelcref{lemitem:tame galois sequence}.
We first note that $K'K^\sh$ agrees with the strict henselisation of $K'$  since $K'K^\sh$ is henselian with separably closed specialisation field and unramified over~$K'$.
Hence, the induced map of the unramified (strongly \'etale) quotients yields an isomorphism onto an open subgroup.
\[
\Gal(K'^\sh/K') \isomto \Gal(K^\sh/K'\cap K^\sh) \subseteq \Gal(K^\sh/K).
\]

For the comparison of tame inertia groups, we note that the inclusion of value groups $\Gamma_K \inj \Gamma_{K'}$
is of finite index $e_{K'/K} \mid [K':K]$. It follows that $\Ext^1(\Gamma_{K'}/\Gamma_K, \hZ'(1))$ is finite and so the induced map 
\[
\Hom(\Gamma_K,\hZ'(1)) \la \Hom(\Gamma_{K'},\hZ'(1))
\]
is injective with open image of finite index. The claim now follows from  the diagram
\[
 \begin{tikzcd}
    1 \ar[r]    & \Hom(\Gamma_{K'},\hZ'(1)) \ar[r] \ar[d,hookrightarrow]   & \Gal(K'^t/K') \ar[r]    \ar[d]  & \Gal(K'^\sh/K') \ar[r]  \ar[d,hookrightarrow]   & 1 \\
    1 \ar[r]    & \Hom(\Gamma_K,\hZ'(1)) \ar[r]   & \Gal(K^t/K) \ar[r]    & \Gal(K^\sh/K) \ar[r]  & 1. 
 \end{tikzcd}
\]
and an application of the snake lemma.
\end{proof}

Also from Kummer theory we obtain Abhyankar's lemma for tame extensions of valued fields:

\begin{lem} \label{lem:tameness-conditions-galois}
    Let $(L,L^+)/(K,K^+)$ be a finite Galois extension of residue characteristic $p \ge 0$.
    The following conditions are equivalent:
    \begin{enumerate}[(a)]
        \item 
        \label{lemitem:tame}
        The extension $(K,K^+) \to (L,L^+)$ is tame.
        \item
        \label{lemitem:tame-galois}
        The wild inertia group~$I_p$ is trivial and the inertia group~$I$ is abelian and via the Kummer pairing isomorphic to $\Hom(\Gamma_L/\Gamma_K,\bQ/\bZ'(1))$.
        \item
        \label{lemitem:tame-abhyankar}
        There exist a finite unramified extension $(K',K'^+)/(K,K^+)$, elements $a_1, \ldots, a_r\in K'$, and an integer $n\geq 1$ not divisible by the residue characteristic such that we have an inclusion of extensions of $K$
        \[ 
            L \subseteq K'(\sqrt[n]{a_1}, \ldots, \sqrt[n]{a_r}).
        \]
    \end{enumerate}
\end{lem}

\begin{proof}
    \labelcref{lemitem:tame} implies \labelcref{lemitem:tame-galois} since $I_p \subseteq I$ is a $p$-group and~$I$ has order prime to~$p$ by assumption.
    
    If~\labelcref{lemitem:tame-galois} holds, the homomorphism (\ref{Kummer-hom}) is an isomorphism as explained above.
    Consequently, the extension $L^\sh/K^\sh$ is abelian and unravelling the definition of the Kummer pairing we see that an element $\gamma \in \Gamma_L/\Gamma_K$ of order~$n$ `corresponds' to the extension $K^\sh(x)/K^\sh$, where $x \in L^\sh$ is such that $\lvert x \rvert = \gamma$ and $x^n \in K^\sh$.
    In total we get that $L^\sh = K^\sh(a_1^{1/q_1},\ldots,a_k^{1/q_r})$  for positive integers~$q_i$ prime to~$p$ and $a_i \in K^\sh$ (see \cite[Corollary 6.2.14]{GabberRamero} for a detailed argument).
    Taking for~$n$ the least common multiple of the~$q_i$, we obtain inclusions
    \[
      L \subseteq L^\sh \subseteq K^\sh(\sqrt[n]{a_1}, \ldots, \sqrt[n]{a_r}).      
    \]
    Since~$L/K$ is finite, $L$ is already contained in $K'(\sqrt[n]{a_1}, \ldots, \sqrt[n]{a_r})$ for a finite subextension~$K'$ of $K^\sh/K$ with $a_i \in K'$.
    This shows the claim of \labelcref{lemitem:tame-abhyankar}.
   
    Finally, from~\labelcref{lemitem:tame-abhyankar} we conclude
    \[
     L^\sh \subseteq K'(\sqrt[n]{a_1}, \ldots, \sqrt[n]{a_r})^\sh = K^\sh(\sqrt[n]{a_1}, \ldots, \sqrt[n]{a_r})
    \]
    and the latter field is abelian of degree prime to~$p$ over~$K^\sh$.
    This shows~\labelcref{lemitem:tame}.
\end{proof}

For extensions that are not necessarily Galois we deduce the following characterisations.

\begin{lem} \label{lem:tameness-conditions}
    Let $(K,K^+) \to (L,L^+)$ be a finite extension of valued fields of residue characteristic~$p$. 
    The following conditions are equivalent: 
    \begin{enumerate}[(a)]
        \item 
        \label{lemitem:tame1}
        The extension $(K,K^+) \to (L,L^+)$ is tame.
        \item
        \label{lemitem:tame-abhyankar1}
        There exist a finite unramified extension $(K',K'^+)/(K,K^+)$, elements $a_1, \ldots, a_r\in K'$, and an integer $n\geq 1$ not divisible by the residue characteristic such that we have an inclusion of extensions of $K$
        \[ 
            L \subseteq K'(\sqrt[n]{a_1}, \ldots, \sqrt[n]{a_r}).
        \]
        \item 
        \label{lemitem:tame-numerical}
        The ramification index $e_{L/K} = \lvert \Gamma_L/\Gamma_K \rvert$ is not divisible by~$p$, the specialisation field extension $L^\succ/K^\succ$ is separable (of degree $f_{L/K}$), and the extension is defectless: the degree $[(L,L^+)^{\rh}:(K,K^+)^{\rh}]$ of the extension of the respectively henselised fields agrees with $e_{L/K} \cdot f_{L/K}$.
    \end{enumerate}
\end{lem}

\begin{proof}
    If \labelcref{lemitem:tame1} holds, the extension $L^\sh/K^\sh$ is abelian by \cref{lem:tameness-conditions-galois}~\labelcref{lemitem:tame-galois}.
    Applying the characterisation 
    \cref{lem:tameness-conditions-galois}~\labelcref{lemitem:tame-abhyankar} to $L^\sh/K^\sh$, we get $a_1,\ldots,a_r \in K^\sh$ and an 
    $n \ge 1$ such that
    \[
        L^\sh \subseteq K^\sh(\sqrt[n]{a_1}, \ldots, \sqrt[n]{a_r}).
    \]
    Choosing a suitable finite subextension $K'$ of $K^\sh/K$ as in the proof of \cref{lem:tameness-conditions-galois} we obtain~\labelcref{lemitem:tame-abhyankar1}. 

    From~\labelcref{lemitem:tame-abhyankar1} we deduce that~$e_{L/K}$ divides $n^r$, which is prime to~$p$.
    Regarding residue field extensions, we know that
    \[
        L^\succ \subseteq L^{\sh\succ} \subseteq K^\sh(\sqrt[n]{a_1}, \ldots, \sqrt[n]{a_r})^\succ = K^{\sh\succ} = K^{\succ\sep}.
    \]
    Moreover, $L^\rh$ is contained in $K'^\rh(\sqrt[n]{a_1}, \ldots, \sqrt[n]{a_r})$, which is a defectless extension of~$K^\rh$, so~$L^\rh$ is also defectless.

    If~\labelcref{lemitem:tame-numerical} holds, then $[L^\sh:K^\sh] = e_{L/K}$ is prime to~$p$ and hence $L/K$ is tame, i.e.\ \labelcref{lemitem:tame1} holds.
\end{proof}


\subsubsection*{Tameness of composite valuations}

In \cref{ss:tameness-comparison-pre} we will need to investigate tame covers of the universal compactification of a rigid space.
The relevant valuations are naturally compositions of two valuations.
As a preparation we study tame extensions of composite valuations.
Our aim is to express the tameness condition in terms of tameness of the composite parts.


\begin{defi}
    Two valued fields $(K_1,K_1^+)$ and $(K_2,K_2^+)$ are said to be \textbf{composable} if $K_1^\succ = K_2$.
    In this case their \textbf{composition} is the valued field $(K,K^+)$ where $K = K_1$ and $K^+$ is the preimage of~$K_2^+$ under the projection $K_1^+ \surj K_1^\succ = K_2$.
    We use the following notation for the composition:
    \[
        (K,K^+) = (K_1,K_1^+) \succ (K_2,K_2^+).
    \]
\end{defi}

The ring $K^+$ in the above definition is indeed a valuation ring of $K$ because for all $x \in K^\times$ we have $x \in K^+$ or $x^{-1} \in K^+$.
The composition $(K,K^+) = (K_1,K_1^+) \succ (K_2,K_2^\succ)$ is illustrated by the following cartesian diagram.

\[
     \begin{tikzcd}
         K = K_1 & K_1^+ \ar[d,twoheadrightarrow] \ar[l,hook']  & K^+ \ar[d,twoheadrightarrow] \ar[l,hook'] \\
         & K_1^\succ = K_2 & K_2^+ \ar[l,hook'] \ar[d,twoheadrightarrow] \\
         & & K_2^\succ = K^\succ.
     \end{tikzcd}
\]

\begin{prop} \label{henselisation-composition}
    Let $(K,K^+) = (K_1,K_1^+) \succ (K_2,K_2^+)$ be a composition of valued fields such that~$K^\succ = K_2^\succ$ is of characteristic $p \ge 0$.
    \begin{enumerate}[ref=(\arabic*)]
        \item \label{lemitem:hens-comp1}
        $(K,K^+)$ is henselian if and only if $(K_1,K_1^+)$ and $(K_2,K_2^+)$ are henselian.
        \item \label{lemitem:hens-comp2}
        $(K,K^+)$ is strictly henselian if and only if $(K_1,K_1^+)$ is henselian and $(K_2,K_2^+)$ is strictly henselian.
        
        \item \label{lemitem:hens-comp3}
        $(K,K^+)$ is tamely henselian 
        if and only if $(K_1,K_1^+)$ is henselian with value group divisible by all primes except possibly~$p$ and $(K_2,K_2^+)$ is tamely henselian.
    \end{enumerate}
\end{prop}

\begin{proof}
    Assertion~\labelcref{lemitem:hens-comp1} has been shown in \cite[Chapter~F, Propositions~9 and 10]{Ribenboim:Theorie-valuations}. 
    From this~\labelcref{lemitem:hens-comp2} follows immediately as a henselian local ring is strictly henselian if and only if its residue field is separably closed and the residue fields of $K^+$ and~$K_2^+$ are the same.
    
    For~\labelcref{lemitem:hens-comp3} remember that a valued field of residue characteristic $p > 0$ is tamely henselian if and only if it is strictly henselian and its value group is divisible by all primes except~$p$.
    The value group $\Gamma_{K_2}$ is a convex subgroup of~$\Gamma_K$ with quotient group~$\Gamma_{K_1}$.
    Therefore, $\Gamma_K$ is divisible by a prime $\ell \ne p$ if and only if $\Gamma_{K_1}$ and~$\Gamma_{K_2}$ are.
    This shows the equivalence.
\end{proof}

\begin{prop} \label{composition-tame}
    Let $(L,L^+)/(K,K^+)$ be an algebraic extension of valued fields of residue characteristic $p \ge 0$.
    Suppose that $(L,L^+) = (L_1,L_1^+) \succ (L_2,L_2^+)$ and set $(K_1,K_1^+) := (K,L_1^+ \cap K)$ and $(K_2,K_2^+) := (K_1^\succ,L_2^+ \cap K_1^\succ)$ so that $(K,K^+) = (K_1,K_1^+) \succ (K_2,K_2^\succ)$.
    Then
    \begin{enumerate}[ref=(\arabic*)]
        \item \label{lemitem:composition etale}
        $(L,L^+)/(K,K^+)$ is unramified if and only if $(L_1,L_1^+)/(K_1,K_1^+)$ and $(L_2,L_2^+)/(K_2,K_2^+)$ are unramified.
        
        In particular, the henselisation and strict henselisation of $(K,K^+)$ decompose as
        \begin{align*}
            (K^\rh,K^{\rh,+}) &= (K'_1,{K'_1}^+) \succ (K_2^\rh,K_2^{\rh,+}),   \\
            (K^\sh,K^{\sh,+}) &= (K''_1,{K''_1}^+) \succ (K_2^\sh,K_2^{\sh,+}),
        \end{align*}
        where $(K'_1,{K'_1}^+)$ and $(K''_1,{K''_1}^+)$ are the unique unramified extensions of $(K_1^\rh,K_1^{\rh,+})$ with residue fields $K_2^\rh$ and $K_2^\sh$, respectively.
        \item \label{lemitem:composition tame}
        $(L,L^+)/(K,K^+)$ is tame if and only if $(L_1,L_1^+)/(K_1,K_1^+)$ and $(L_2,L_2^+)/(K_2,K_2^+)$ are tame and the ramification index of $(L_1,L_1^+)/(K_1,K_1^+)$ is prime to~$p$ \footnote{It might happen that the residue characteristic of $(K_1,K_1^+)$ is~$0$ and the residue characteristic of $(K_2,K_2^+)$ is $p > 0$.
        In this case extracting $p$-th roots in $(K_1,K_1^+)$ is tame but does not lead to a tame extension of $(K,K^+)$.
        This is the reason for the extra condition on the ramification index.}.

        In particular, the tame henselisation of $(K,K^+)$ decomposes as
        \[
            (K^t,K^{t,+}) = (K'_1,{K'_1}^+) \succ (K_2^t,K_2^{t,+}),
        \]
         where $(K'_1,{K'_1}^+) \subseteq (K_1^t,K_1^{t,+})$ is the minimal extension of $(K_1^\rh,K_1^{\rh,+})$ with residue field~$K_2^t$ and value group divisible by all primes except possibly~$p$.
    \end{enumerate}
\end{prop}

\begin{proof}
    Let $\fp \subseteq K^+$ be the maximal ideal of~$K_1^+$, so that $K_1^+ = K^+_\fp$ and $K_2^+ = K^+/\fp$.
    
    In order to prove~\labelcref{lemitem:composition etale} suppose that $(L,L^+)/(K,K^+)$ is unramified.
    Then $\fp_L := \fp L^+$ is the maximal ideal of~$L_1^+$.
    Moreover, $L_1^+$ and $L_2^+$ are the base changes of~$L^+$ via $K^+ \hookrightarrow K_1^+$ and $K^+ \twoheadrightarrow K_2^+$, respectively.
    In particular, both extensions $L_1^+/K_1^+$ and $L_2^+/K_2^+$ are unramified.
    
    Suppose now that $(L_1,L_1^+)/(K_1,K_1^+)$ and $(L_2,L_2^+)/(K_2,K_2^+)$ are unramified.
    We can check that $(L,L^+)/(K,K^+)$ is unramified after base changing to $(K^\sh,K^{\sh,+})$.
    So we may assume that $(K,K^+)$ is strictly henselian.
    By \cref{henselisation-composition}~\labelcref{lemitem:hens-comp2}, also $(K_2,K_2^+)$ is strictly henselian and $(K_1,K_1^+)$ is henselian.
    Therefore, the unramified extension $(L_2,L_2^+)/(K_2,K_2^+)$ is trivial.
    Now $(L_1,L_1^+)/(K_1,K_1^+)$ is an unramified extension with trivial residue field extension of a henselian valued field, hence trivial as well.
    This shows that $(L,L^+) = (K,K^+)$, which is trivially unramified.

    For the second part of~\labelcref{lemitem:composition etale} we note that the unramified extension $(K^h,K^{h,+})/(K,K^+)$ induces unramified extensions of $(K_1,K_1^+)$ and $(K_2,K_2^+)$ that have to contain the respective henselisations by \cref{henselisation-composition}~\labelcref{lemitem:hens-comp1}.
    Invoking that the residue field extension over~$K_2^\succ$ has to be trivial we obtain the desired decomposition.
    The argument for the strict henselisation is analogous.

    In order to prove~\labelcref{lemitem:composition tame} we first assume that $(L,L^+)/(K,K^+)$ is tame.
    We already know from~\labelcref{lemitem:composition etale} that passing to the strict henselisation of $(K,K^+)$ induces unramified extensions of $(K_1,K_1^+)$ and $(K_2,K_2^+)$.
    We may therefore assume that $(K,K^+)$ is strictly henselian.
    Then $L/K$ is of degree prime to~$p$ and the same holds for $L_1/K_1$ and $L_2/K_2$, so both extensions $(L_1,L_1^+)/(K_1,K_1^+)$ and $(L_2,L_2^+)/(K_2,K_2^+)$ are tame and the ramification index of the former extension is prime to~$p$.
    
    For the converse we reduce to the situation where $(K,K^+)$ is tamely henselian.
    Then we argue in the same way as before using that a tame extension of henselian valued fields with trivial residue field extension and trivial value group extension is trivial.
\end{proof}


\begin{rmk}
    The extension $\bQ_p((T^{1/p}))/\bQ_p((T))$ certainly defines a tame extension of $(\bQ_p((T)),\bQ_p[[T]])$ since the residue field has characteristic~$0$.
    However, if we consider it over the composition
    \[
     (\bQ_p((T)),\bQ_p[[T]]) \succ (\bQ_p,\bZ_p) = (\bQ_p((T)),\bZ_p + T\bQ_p[[T]]),
    \]
    the same extension is not tame as it is ramified of degree~$p$ and the residue characteristic is now~$p$.
    One might be tempted to declare such an extension to be tame since it induces tame extensions over both constituents $(\bQ_p((t)),\bQ_p[[T]])$ and $(\bQ_p,\bZ_p)$.
    However, this is not a good idea for various reasons.

    Firstly, the module of log differentials for the extension over $(\bQ_p((T)),\bZ_p + T\bQ_p[[T]])$ is generated by $\dlog T^{1/p}$ subject to the relation $0 = \dlog T = p \dlog T^{1/p}$.
    So it is isomorphic to $\bZ/p\bZ$, which is not zero as expected for tame extensions.

    Secondly, the base change of the extension to $\bQ_p(((pT)^{1/p}))$ transforms it into the extension $\bQ_p((S))(p^{1/p})/\bQ_p((S))$, where $S = (pT)^{1/p}$.
    This extension is certainly wild with respect to the valuation ring $\bZ_p + T\bQ_p[[S]]$ but a base change of a tame extension should be tame.

    In what follows, we will work with rigid spaces in mixed or positive characteristic.
    In this setup the above described subtlety does not come up as for all relevant valued fields $(K,K^+)$ the residue characteristic of $K^\circ$ (the power bounded elements) is~$p$.
    In other words: the only possible localisation of~$K^+$ of residue characteristic~$0$ equals~$K$.
\end{rmk}


%%%---------------------------------------
\subsection{Tame covers of schemes} 
\label{sec:adic tame discrete case}
%%%---------------------------------------

In this section, we introduce tame covers of schemes over a base scheme~$S$.
Given its valuation-theoretic nature, the tameness of an \'etale (or just locally quasi-finite) morphism of schemes $Y\to X$ relative to $S$ is best defined using the associated map of discretely ringed adic spaces $\Spa(Y, S)\to \Spa(X, S)$, see \cite[\S3.1]{Temkin2011:RelativeRiemannZariskiSpacesa} and \cite{Hubner2021:AdicTameSite}.

In the literature, tame covers of schemes have been studied for decades, but there are varying notions of tameness used for their definition. 
After giving the definition of tame covers we use in this article (sometimes we refer to this as \emph{adic tameness} because it relies on tameness of the associated morphism of adic spaces), we compare it to the existing notions of tameness.
We will see that under some mild assumptions most of the tameness notions are equivalent.
This was already established in \cite{KerzSchmidt2010:DifferentNotionsTameness}.
However, they had to assume resolution of singularities for some of the implications and we will report on new developments that allow for a proof without using resolutions.

\begin{defi}
\label{defi:Spa for schemes}
Let $f\colon X\to S$ be a morphism of schemes. We define $\Spa(X,S)$ as the set of triples $(x,V,s)$, where $x \in X$ is a point, $V$ is a valuation subring of $k(x)$, and $s \colon \Spec(V) \to S$ is a map (that we call a \textbf{centre of~$V$ on~$S$}) fitting inside a commutative square
\[
 \begin{tikzcd}
     \Spec(k(x)) \ar[r] \ar[d] & X \ar[d] \\
     \Spec(V) \ar[r,"s"] & S.
 \end{tikzcd}
\]
\end{defi}

Often we denote a triple $(x,V,s)\in \Spa(X,S)$ simply by $x$, and the valuation ring $V\subseteq k(x)$ will be denoted by $k(x)^+$. By abuse of notation we sometimes also write~$s$ for the image of the closed point under the map $\Spec(V) \to S$ in the definition.
This is justified as for a given point $s \in S$ there is at most one map $\Spec(V) \to S$ as above mapping the closed point of~$\Spec(V)$ to~$s$.

    A commutative square 
    \[ 
        \begin{tikzcd}
            X'\ar[d] \ar[r,"g"] & X\ar[d] \\
            S'\ar[r,"h"] & S
        \end{tikzcd}
    \]
    induces a map $\Spa(X', S')\to \Spa(X, S)$ sending $(x', V',s')$ to $(x,V,s)$ where $x=g(x')$, $V = V'\cap k(x)$, and $s$ is induced by $h\circ s'$.

The construction of $\Spa(X,S)$ was introduced in \cite[\S~3.1]{Temkin2011:RelativeRiemannZariskiSpacesa}, where it is also shown that $\Spa(X,S)$ naturally carries the structure of an adic space.
In \cref{s:compactifications} we will see that $\Spa(X,S)$ is a special instance of a vertical compactification, namely of $\Spa(X,X)$ over $\Spa(S,S)$ (see \cref{example-Spa-schemes}).


\begin{defi}
\label{defi:tame cover of schemes X over S}
    Let $f\colon Y \to X$ be an \'etale morphism of schemes over $S$.   
    \begin{enumerate}
        \item 
        We say $f\colon Y \to X$ is \textbf{tame relative to $S$} if for every point $y = (k,k^+)$ in $\Spa(Y, S)$ with image $x = f(y)$ in $\Spa(X, S)$ the induced extension
        $
            \big(k(y),k(y)^+\big) / \big(k(x), k(x)^+\big)
        $
        is tamely ramified.
        \item 
        The category of all finite \'etale morphisms $Y\to X$ which are tame relative to $S$ is denoted by $\FEtt_{X/S}\subseteq \FEt_X$. 
    \end{enumerate}
\end{defi}

The base change functor $g^\ast \colon \FEt_X \to \FEt_{X'}$ induces a base change functor
\[
    g^\ast \colon \FEtt_{X/S} \la \FEtt_{X'/S'}
\]
on the respective full subcategories of tame covers relative to $S$ (resp. to $S'$).


%%%---------------------------------------
\subsubsection*{Comparison with other notions of tameness}

In \cite{KerzSchmidt2010:DifferentNotionsTameness} Kerz and Schmidt compare various tameness conditions that we review quickly in order to compare with (adic) tameness as understood in this note. 
Let $X$ be a separated scheme of finite type over an integral, pure-dimensional, separated and excellent base scheme $S$.
Then we consider for a finite \'etale map $Y \to X$ the following properties. 

The oldest definition of a tame cover goes back to \cite[Exposé~X]{SGA1} and requires~$X$ to be regular with a regular compactification~$\ov X$, proper over the base~$S$, whose boundary is the support of a normal crossing divisor $D$.
In this setting we define \textbf{Grothendieck-Murre tameness} as follows:
A finite étale cover $Y \to X$ is \textbf{tame along~$D$} if the normalisation $\ov{Y} \to \ov{X}$ is tamely ramified at rank~$1$ discrete valuation rings corresponding to irreducible components of $D$.

\begin{prop}
    For~$X$ and~$S$ as above, a finite étale cover $Y \to X$ is tame relative to~$S$ if and only if it is tame along~$D$.
\end{prop}

\begin{proof}
    By the valuative criterion of properness we may assume that $S = \ov X$.
    A discrete valuation ring~$V_P$ corresponding to the irreducible component~$P$ of~$D$ defines the point $(\eta,V_P,\eta_P)$ of $\Spa(X,\ov X)$, where~$\eta$ is the generic point of~$X$ and~$\eta_P$ the generic point of~$P$.
    Tameness of $Y \to X$ along~$D$ translates to tameness of $\Spa(Y,\ov X) \to \Spa(X,\ov X)$ over these points $(\eta,V_P,\eta_P)$.
    In particular, tameness of $Y \to X$ relative to~$S$ implies tameness along~$D$.

    In order to show that Grothendieck-Murre tameness along~$D$ implies tameness relative to~$\ov X$ we reduce to the case where~$\ov X = \Spec(A)$ is strictly henselian.
    Then Abhyankar's lemma \cite[XIII Proposition~5.2]{SGA1} tells us that~$Y$ is dominated by a cover of the form
    \[
     \Spec(A[a_1^{1/n},\ldots,a_r^{1/n}]) \longrightarrow \Spec(A),
    \]
    where $a_1,\ldots,a_r \in A$ define the irreducible components of~$D$ and~$n$ is prime to the residue characteristic~$p$ of~$A$.
    Since the degree of this cover is prime to~$p$, it is tame (compare with the proof of \cite[Proposition~4.2]{KerzSchmidt2010:DifferentNotionsTameness}).
\end{proof}

We could rephrase the proof above as Abhyankar's lemma telling us that $\ov Y \to \ov X$ is Kummer étale with respect to the log structure defined by~$D$.
We will encounter a variant of this argument in a more sophisticated (nonnoetherian) setting in \cref{ss:log-abhyankar}.

Grothendieck-Murre tameness is a well-behaved concept only in the presence of a normal crossing compactification.
In order to define tame covers also under less favourable circumstances, various notions of tameness have been introduced.
For a separated scheme~$X$ of finite type over the base~$S$ (which is still assumed to be integral, pure-dimensional, separated, and excellent), we have the following list of tameness conditions for a finite étale cover $Y \to X$:
\begin{enumerate}
    \item \label{def:curve tameness}
    \textbf{Curve tameness} (following Wiesend \cite{Wiesend2008:TamelyRamifiedCovers}):    
    the pull-back $Y_C \to C$ along $C \to X$ to every connected scheme $C$ of finite type over $S$ with $\dim_S C = 1$ (see \cite[Definition on p.653]{KerzSchmidt2010:DifferentNotionsTameness}, this might deviate from the Krull dimension) is tame.
    Any such $C$ has a canonical set of discrete valuations of its function field as points at infinity with respect to which tameness can be defined.
\end{enumerate}
If in addition~$X$ is normal, the following tameness concepts are available:
\begin{enumerate}
    \setcounter{enumi}{1}
    \item  \label{def:divisor tameness}
    \textbf{Divisor tameness:} every discrete rank $1$ valuation of the function field $k(X)$ of $X$ arising from a boundary divisor in a normal compactification $X \subseteq \ov X$ over $S$ is tamely ramified in the function field of $Y$.
    \item  \label{def:chain tameness}
    \textbf{Chain tameness:} for every normal compactification $\ov X \to S$ of $X$,
    every discrete rank $d = \dim_S(X)$ valuation of $k(X)$ dominating a Parshin chain on $\ov X$ whose valuation ring dominates a local ring $\cO_{S,s}$ on $S$ is tamely ramified in $k(Y)$.
    \item  \label{def:valuation tameness}
    \textbf{Valuation tameness:} every valuation of $k(X)$ with centre on $S$ is tamely ramified in $k(Y)$.   
\end{enumerate}
All of the above definitions can be interpreted as placing tameness conditions on points of $\Spa(Y,S)$ of a special kind and are thus implied by adic tameness.
For instance, for curve tameness we only use points $(y,V,s) \in \Spa(Y,S)$ where~$y$ is the image of the generic point of a test curve~$C$ as in the definition above.
For curve tameness, divisor tameness and chain tameness these special points arise from the geometry of the scheme~$Y$.
The corresponding valuations are rather well behaved and geometric in nature.
In \cite{KerzSchmidt2010:DifferentNotionsTameness} the authors show that these notions of tameness are equivalent using some smart scheme theoretic arguments to obtain the following result.

\begin{prop}[\cite{KerzSchmidt2010:DifferentNotionsTameness}, Theorem~4.4]
    Let~$X$ be a regular, separated scheme of finite type over~$S$.
    Then curve tameness, divisor tameness, and chain tameness are equivalent for finite étale covers of~$X$.
\end{prop}

Valuation tameness includes tameness conditions for arbitrary $S$-valuations of the function field of~$X$.
It is a priori stronger than the above more geometric notions of tameness and closer to our (adic) definition of tameness.
Accordingly, it is not so hard to prove a comparison between the two as follows.

\begin{prop}[\cite{Hubner2021:AdicTameSite}, Proposition~9.4]
    Let $X/S$ be regular, separated, and of finite type.
    Then adic tameness of finite étale covers of~$X$ is equivalent to valuation tameness.
\end{prop}

In \cite{KerzSchmidt2010:DifferentNotionsTameness} Kerz and Schmidt bridge the gap between valuation tameness and the more geometric tameness notions by resorting to resolution of singularities.
Using alterations, Schmidt proves an unconditional comparison for regular varieties over a field of positive characteristic  in \cite[Theorem~4.4]{Schmidt-quasi-purity}.
The following is the most general comparison theorem, accomplished by using non-archimedean geometry in order to avoid the need for resolution of singularities.

\begin{prop}[\cite{HuebnerTemkin2026:WildLocus}, Theorem~8.11, Corollary~8.12]
    Suppose~$S$ is of finite type over a Dedekind scheme and~$X$ is separated and of finite type over~$S$.
    \begin{enumerate}
        \item If~$X$ is regular, valuation tameness for finite étale covers over~$X$ implies adic tameness.
        \item Curve tameness for finite étale covers over~$X$ implies adic tameness.
    \end{enumerate}
\end{prop}

The last notion of tameness we would like to discuss is \textbf{numerical tameness} with respect to a fixed normal compactification $\ov X \to S$ of $X$:
a finite étale cover $Y \to X$ is numerically tame if for all geometric points $\bar x \in \ov{X} \setminus X$ the restriction of $Y \to X$ to the scheme of nearby points (also called the link of $X$ near $\bar x$), $ U_{\bar x}^{\sh} = X \times_{\bar X} \Spec(\cO^\sh_{\bar X, \bar x}) $ decomposes into connected components whose Galois closures have degree prime to the characteristic of $\bar x$.

This notion of tameness differs from the preceding ones in two aspects: it requires a fixed normal compactification and it does not involve valuations.
In fact, numerical tameness depends on the choice of the compactification as Examples~2 and~3 in the appendix of \cite{KerzSchmidt2010:DifferentNotionsTameness} illustrate.
It is a priori stronger than the above discussed notions of tameness.

\begin{prop}[\cite{KerzSchmidt2010:DifferentNotionsTameness}, Theorem~5.3] \label{numerical-implies-tame}
    Let~$X$ be normal, separated, and of finite type over~$S$ with a normal compactification~$\ov X$.
    Then numerical tameness of finite étale covers of~$X$ along $\ov X \setminus X$ implies adic tameness (and hence curve, divisor, and chain tameness).
\end{prop}

\begin{proof}
    A variant of this has been shown in \cite[Theorem~5.3]{KerzSchmidt2010:DifferentNotionsTameness} for regular schemes~$X$ and divisor tameness as a conclusion.
    A similar argument can be used here:
    Let $Y \to X$ be a finite étale cover and denote by $\ov Y$ the normalisation of~$\ov X$ in~$Y$.
    We start with a point $(x,V,s)$ of $\Spa(X,S)$ and denote by $\bar{x}$ the unique centre of~$x$ on~$\ov X$ lying over~$s$.
    The centre map $\Spec(V^\sh )\to \ov X$ naturally factors through $\Spec(\cO_{\ov X,\bar{x}}^\sh)$.
    Now by definition the cover $Y \to X$ restricted to the link~$U_{\bar x}^{\sh}$ splits into connected components with Galois closures of degree prime to the residue characteristic~$p$ of~$\bar{x}$.
    Restricting further to $X \times_{\ov X} \Spec(V^\sh)$ this remains true, which implies that $Y \to X$ is tame at all points of $\Spa(Y,S)$ over $(x,V,s)$.
\end{proof}

A converse holds if we are allowed to change the compactification.

\begin{prop}[\cite{HuebnerTemkin2026:WildLocus}, Theorem~9.3]
    Suppose that~$X$ is normal, separated, and of finite type over~$S$.
    If $Y \to X$ is a finite étale cover that is tame, there exists a compactification~$\ov X$ of~$X$ over~$S$ such that $Y \to X$ is numerically tame along $\ov X \setminus X$.
\end{prop}

The following diagram summarises the notions of tameness we have discussed.

\[
\begin{tikzcd}[column sep=7em, row sep=4em]
\fbox{\begin{tabular}{c}
numerical \\
tame wrt $X \subseteq \ov{X}$
\end{tabular}} 
\ar[r, Rightarrow, shift left=1.5ex, "\text{Prop}~\ref{numerical-implies-tame}"] 
& 
\fbox{\begin{tabular}{c}
(adic) tame \\
relative $S$
\end{tabular}}
\ar[l, Rightarrow, shift left=1.5ex, "\substack{\text{[HT26, Thm~9.3]} \\ \exists \ov X}"]
\ar[d, Rightarrow, shift left=1.5ex, "\text{def}"] 
\ar[r, Rightarrow, shift left=1.5ex, "\text{def}"] 
& 
\fbox{\begin{tabular}{c}
curve \\
tame
\end{tabular}}
\ar[l, Rightarrow, shift left=1.5ex, "\text{[HT26, Cor~8.12]}"]
\ar[d, Rightarrow, shift left=1.5ex, "\substack{\text{[KS10,} \\ \text{Lem~2.4]}}"]
\\
& 
\fbox{\begin{tabular}{c}
valuation \\
tame
\end{tabular}}
\ar[u, Rightarrow, shift left=1.5ex, "\substack{\text{[HT26,} \\ \text{Thm~8.11]} \\ \text{$X$ regular}}"]
\ar[r, Rightarrow, shift left=1.5ex, "\text{def}"]
&
\fbox{\begin{tabular}{c}
divisor \\
tame
\end{tabular}}
\ar[l, Rightarrow, shift left=1.5ex, "\text{[HT26, Thm~8.11]}"]
\ar[u, Rightarrow, shift left=1.5ex, "\substack{\text{[KS10,} \\ \text{Thm~4.4]} \\ \text{$X$ regular}}"] 
\ar[d, Rightarrow, shift left=1.5ex, "\substack{\text{[KS10,} \\ \text{Lem~3.5]}}"]
\\
&
&
\fbox{\begin{tabular}{c}
chain \\
tame
\end{tabular}}
\ar[u, Rightarrow, shift left=1.5ex, "\substack{\text{[KS10,} \\ \text{Thm~4.4]} \\ \text{$X$ regular}}"]
\end{tikzcd}
\]

%%%---------------------------------------
\subsection{Review of Galois categories}
\label{sec:Gal cats}
%%%---------------------------------------

We recall the definition of a Galois category from 
\stacks[Definition]{0BMY}.  These axioms are stronger than  but ultimately equivalent to the original axioms from \cite[V, \S4]{SGA1}. In particular, we use the term \emph{connected object} in the sense of \stacks{0BMY}: an object $X$ is \textbf{connected}  if it is not the initial object and every monomorphism $Y \to X$ is an isomorphism or $Y$ is the initial object. 

Unlike \stacks{0BMY}, we do not consider a fixed fibre functor to be part of the datum of a Galois category.

\begin{defi}
    \label{defi:Gal cat}
    A \textbf{Galois category} is a category $\cC$ that admits a \textbf{fibre functor} $F \colon \cC \to \Sets$ such that the following axioms hold. 
    \begin{enumerate}[(G\arabic*)]
        \item 
        The category $\cC$ has finite limits and finite colimits.
        \item 
        Every object of $\cC$ is a finite (possibly empty) coproduct of connected objects. 
        \item 
        The functor $F$ takes values in finite sets.
        \item 
        The functor $F$ is exact 
        and conservative. 
    \end{enumerate}
    A \textbf{morphism} between Galois categories is an exact functor.    
    Recall that a functor $F$
    \begin{itemize}
        \item 
        is called \textbf{exact} if it commutes with finite limits and finite colimits, and 
        \item 
        it is \textbf{conservative} if a morphism $f$ is an isomorphism if and only if $F(f)$ is an isomorphism. 
    \end{itemize} 
\end{defi}

The main theorem of Galois categories \cite[V~Th\'eor\`eme 4.1]{SGA1}, see also \stacks{0BMR} and \stacks{0BN4}, states the following. The group
\[
    \pi_1(\cC,F) =  \Aut(F \colon \cC \to \Sets),
\]
which is called the \textbf{fundamental group of $\cC$ with base point $F$}, 
is a profinite group (with topology induced by declaring stabilisers of points $x \in F(X)$ to be open). Moreover, the natural enrichment
\[
    F \colon \cC \longrightarrow \pi_1(\cC,F)\text{-}\sets
\]
is an equivalence of $\cC$ with the category of finite sets with continuous action of $\pi_1(\cC,F)$. Furthermore, by \cite[V Proposition~6.1]{SGA1} and similarly \stacks{0BN5}, composition with an exact functor transforms fibre functors into fibre functors. In fact this property of a functor between Galois categories is equivalent to being a morphism.

As in \cref{defi:tame cover of schemes X over S}, we will frequently encounter a full subcategory $\cC'$ of a Galois category $\cC$ and would like to deduce that $\cC'$ is Galois with respect to the restriction of the fibre functor. 

\begin{lem} 
\label{lem:Galois-subcategory-criterion} 
    Let $\cC$ be a Galois category with fibre functor $F \colon \cC\to \mathrm{Set}$, and let $\cC'$ be a full
    subcategory of $\cC$ containing $1_{\cC}$. Then, $\cC'$ is a Galois category with fibre functor $F|_{\cC'}$ if and only if the following conditions are satisfied:
    \begin{enumerate}[label=(\roman*)]
        \item 
        \label{lemitem:fullGalCat1}
        if $X'\to X$ is an epimorphism and $X'$ is isomorphic to an object of $\cC'$, then so is $X$,
        \item 
        \label{lemitem:fullGalCat2}
        an object of $\cC$ is isomorphic to an object of  $\cC'$ if and only if all of its connected components are,
        \item 
        \label{lemitem:fullGalCat3}
        the product of two objects in $\cC'$ is isomorphic to an object of $\cC'$.
    \end{enumerate}
\end{lem}

\begin{proof}
    For simplicity, we may first replace $\cC'$ by the strictly full subcategory of $\cC$ of all objects isomorphic to an object of $\cC'$. Moreover, the fibre functor yields an equivalence of $\cC$ with $\pi\text{-}\sets$, the category of finite sets with continuous $\pi$-action for the profinite group $\pi = \pi_1(\cC,F)$. We may therefore replace $\cC$ by $\pi\text{-}\sets$. 

    If $\cC'$ is Galois, then clearly 
    \labelcref{lemitem:fullGalCat1}--\labelcref{lemitem:fullGalCat3} hold. So we now suppose conditions \labelcref{lemitem:fullGalCat1}--\labelcref{lemitem:fullGalCat3} hold. 
    For every object $X$ of $\cC'$ and any point $x \in F(X)$ the stabiliser in $\pi$ is denoted by $U_{X,x}$. This is an open subgroup of $\pi$ and the $\pi$-set $\pi/U_{X,x}$ is isomorphic to the connected component of $X$ containing $x$. The object $\pi/U_{X,x}$ belongs to $\cC'$ by 
    \labelcref{lemitem:fullGalCat2}. The family of subgroups $U_{X,x}$ is closed under finite intersections by \labelcref{lemitem:fullGalCat3}, 
    under passing to a larger open subgroup by \labelcref{lemitem:fullGalCat1}, and under conjugation by arbitrary elements $\gamma \in \pi$ since $U_{X,\gamma(x)} = \gamma U_{X,x} \gamma^{-1}$.

    The intersection $N$ of all $U_{X,x}$ is a closed normal subgroup of $\pi$, and $\cC'$ is contained in the full subcategory $\pi/N\text{-}\sets$ of $\pi\text{-}\sets$. We claim that $\cC' \simeq \pi/N\text{-}\sets$. Indeed, since by \labelcref{lemitem:fullGalCat2} both are determined by their connected objects it suffices to see that any connected $X$ in $\pi/N\text{-}\sets$ actually belongs to $\cC'$. Such an object is isomorphic to a $\pi/U$ for an open subgroup $U$ that contains $N$. By the definition of $N$ as an intersection of the $U_{X,x}$ and compactness, there are finitely many $X_i$ and $x_i$ such that $U$ contains the intersection of the $U_{X_i,x_i}$. But then, $U$ itself is of the form $U_{X,x}$ for some $X$ and $x$. This shows that $\pi/U$  belongs to $\cC'$.
\end{proof}

\begin{lem}
\label{lem:filtered colimits of Galois categories}
   Let $\cC_i$ be a filtered system of Galois categories such that for $i < j$ the functor $\cC_i \to \cC_j$ is exact. Let $\cC = \varinjlim_i \cC_i$ be the colimit in the $(2,1)$-category of small categories\footnote{This colimit exists by  \cite[3.74]{Kelly:EnrichedCats}.}, and let $F\colon \cC \to \Sets$ be a functor such that each composition $F_i\colon \cC_i \to \cC \to \Sets$ is a fibre functor. 
   
   Then $\cC$ is a Galois category with fibre functor $F$, the functors $\cC_i \to \cC$ are exact, and the natural map 
   \[
   \pi_1(\cC,F) \longrightarrow \varprojlim_i \pi_1(\cC_i, F_i)
   \]
   is an isomorphism of profinite groups.
\end{lem}

\begin{proof}
    The filtered system $(\cC_i,F_i)$ yields a cofiltered system of profinite groups 
    $\pi_i = \pi_1(\cC_i,F_i)$. Using the $F_i$ we may replace $\cC_i$ by $\pi_i$-$\sets$ and $F_i$ by just the forgetful functor into the category of sets. The colimit $\cC$ can then be computed as the category whose objects are finite sets with an action of $\pi_i$ for $i \gg 0$ (equivalently some $i$ or compatibly for all large enough $i$). Morphisms in $\cC$ are given by 
    \[
    \Hom_{\cC}(X,Y) = \varinjlim_i \Hom_{\cC_i}(X,Y)= \varinjlim_i \Hom_{\pi_i}(X,Y).
    \]
    This category agrees with the category of finite sets with continuous $\pi = \varprojlim_i \pi_i$-action, and $F$ enriches to an equivalence $F \colon \cC \isomto \pi\text{-}\sets$.  
\end{proof}


%%%---------------------------------------
\subsection{The tame fundamental group of a scheme}
\label{sec:tame pi1}
%%%---------------------------------------

We come back to the Galois theory of tame covers of an $S$-scheme $X$.

\begin{prop}
\label{prop:Fett is Galois for schemes}
    Let $X$ be a non-empty connected scheme, and let $X\to S$ be a morphism of schemes. Then $\FEtt_{X/S}$ is a Galois category and the inclusion $\FEtt_{X/S} \subseteq \FEt_X$ is exact.
\end{prop}

\begin{proof}
We know from \stacks[Lemma]{0BNB} that $\FEt_X$ is a Galois category. Moreover, the conditions of \cref{lem:Galois-subcategory-criterion} are satisfied by the inclusion $\FEtt_{X/S}\subseteq \FEt_X$. 
\end{proof}

\begin{cor}
\label{cor:exact functor tame covers}
    The functor $g^\ast \colon \FEtt_{X/S} \to \FEtt_{X'/S'}$  
    is an exact functor of Galois categories if $X$ and $X'$ are connected and non-empty.
\end{cor}

\begin{proof}
    This is simply the restriction of the exact functor $g^\ast \colon \FEt_X \to \FEt_{X'}$.
\end{proof}

\begin{defi}
\label{defi:tamepi1 schemes}
    Let $X$ be a connected $S$-scheme, and let $\bar x \to X$ be a geometric point. We denote by $F_{\bar x}$ the fibre functor on $\FEtt_{X/S}$ implicit in \cref{prop:Fett is Galois for schemes} that is the restriction of the fibre functor `fibre at $\bar x$' from $\FEt_X$.     
    \begin{enumerate}
        \item 
        The \textbf{tame fundamental group} of $X$ (relative to $S$) is the profinite fundamental group  
        \[
        \pitame(X/S,\bar x) = \pi_1(\FEtt_{X/S}, F_{\bar x}) = \Aut(F_{\bar x} \colon \FEtt_{X/S} \to \Sets)
        \]
        of the Galois category $\FEtt_{X/S}$. 
        \item
        In view of \cref{cor:exact functor tame covers} a    commutative square
        \[ 
        \begin{tikzcd}
            X'\ar[d] \ar[r,"g"] & X\ar[d] \\
            S'\ar[r] & S
        \end{tikzcd}
        \]
        with $X$ and $X'$ connected together with geometric points $\bar x' \to X'$ and $\bar x = g(\bar x')$ induces in a functorial way a group homomorphism
        \[
        g_\ast \colon \pitame(X'/S',\bar x')  \la \pitame(X/S,\bar x).
        \]
\end{enumerate}
\end{defi}

\begin{rmks}
    \begin{enumerate}
        \item 
        For morphisms $X \to S' \to S$, the fully faithful inclusions of Galois categories $\FEtt_{X/S}  \subseteq \FEtt_{X/S'}  \subseteq\FEt_X$ yield canonical surjective homomorphisms
        \[
        \pi_1(X,\bar x) \surj \pitame(X/S',\bar x) \surj \pitame(X/S,\bar x).
        \]
        These are just the maps $\id_{X,\ast}$, and if $X \to S$ is the 
        identity $\id\colon X \to X$, then the canonical map is an isomorphism
        \[
        \pi_1(X,\bar x) \isomto \pitame(X/X,\bar x).
        \]
     
        \item 
        The tame fundamental group of $X$ relative to $S$ was considered in \cite[\S5]{HubnerSchmidt2021:TameSiteSchemea} in the locally noetherian case and defined as the pro-discrete monodromy group of local systems on the locally connected tame site $\Spa(X,S)_{\rt}$ in the spirit of the \emph{pro-groupe fondamentale \'elargi}, see \cite[X~\S6]{SGA3:2}.  Upon profinite completion, we recover our $\pitame(X/S,\bar x)$, see  also \cite[\S9]{Hubner2021:AdicTameSite}.      
    \end{enumerate}
\end{rmks}

\begin{rmk}
For a connected and separated scheme of finite type~$X$ over an integral, pure-dimensional, and excellent scheme~$S$ and a geometric point $\bar{x} \to X$, the curve tame covers form a Galois category and give rise to the curve tame fundamental group of~$X$.
If, in addition,~$X$ is normal, we can form the chain tame, divisor tame, and valuation tame fundamental groups.
If~$X$ is regular with an snc compactification~$\ov X$, we furthermore have the Grothendieck-Murre tame fundamental group at our disposal, which we can identify with the Kummer étale fundamental group of~$\ov X$ endowed with the compactifying log structure defined by~$X$. 

Under the various hypotheses discussed in \cref{sec:adic tame discrete case} all of the above fundamental groups are isomorphic to the tame fundamental group $\pi_1^t(X/S,\bar{x})$.
We can thus reinterpret properties of some variant of a tame fundamental group proved in the literature as results about the (adic) tame fundamental group.
\end{rmk}


%%%---------------------------------------
\subsection{Descent for Galois categories and graphs of groups}
\label{sec:descent for Gal cats}
%%%---------------------------------------

We review the content of \cite{Stix2006:vanKampen}, a profinite version of \cite{HaefligerComplexesGroups}, with emphasis on the application to finite generation results. 

%%%---------------------------------------
\subsubsection*{Complex of Galois categories}
Let $\Delta_{\leq 2}$ be the category with three objects $[n] = \{0,\ldots,n\}$ for $0 \leq n \leq 2$. Morphisms in $\Delta_{\leq 2}$ are \emph{strictly} monotone maps of sets. We will use the following notation for maps in $\Delta_{\leq 2}$:
\[
\begin{array}{rll}
    v_i \colon & [0] \to [2] & \qquad \text{ image equals $\{i\}$ and $i \in \{0,1,2\}$,} \\[1ex]
    \pr_{i} \colon & [0] \to [1] & \qquad \text{ image equals $\{i\}$ and $i \in \{0,1\}$,} \\[1ex] 
    \pr_{ij} \colon & [1] \to [2] & \qquad  \text{ image equals $\{i,j\}$ and $i,j \in \{0,1,2\}$.}
\end{array}
\]

A \textbf{$2$-complex of sets} is a functor $E \colon \Delta_{\leq 2}^\op  \to \Sets$, 
i.e.~essentially a diagram in the category of sets of the form
\[
E_\bullet = \Big[\begin{tikzcd}
E_0 & \ar[l,shift right=0.7ex] \ar[l,shift left=0.7ex] E_1 & \ar[l,shift right=1ex] \ar[l,shift left=1ex] \ar[l] E_2  
\end{tikzcd}\Big]
\]
with the usual simplicial identities among compositions of maps. 
We also consider $E_\bullet$ as a category: objects of $E_\bullet$ are the elements of the sets $E_i$ and morphisms are built from $\partial\colon [i] \to [j]$ in $\Delta_{\leq 2}$ and $t \in E_j$ with $s = E(\partial)(t)$ as  morphisms $\partial^\ast\colon s \to t$. Composition in $E_\bullet$ is derived from composition in $\Delta_{\leq 2}$. 

\begin{defi}
    We say that $E_\bullet$ is \textbf{connected} if the topological realisation $|E_\bullet|$ is connected. 
\end{defi}

Let $E_\bullet$ be a $2$-complex of sets. Let $\cG_\bullet \to E_\bullet$ be a fibred category in Galois categories, i.e. for all $s \in E_i$ the category $\cG_s = \cG_\bullet(s)$ is a Galois category and the pull-back functors are exact. The associated $2$-complex of multi-Galois categories (see \cite[V~\S9]{SGA1})
also denoted by~$\cG_\bullet$ is 
\[
\cG_\bullet = \left[\begin{tikzcd}
\cG_0 \ar[r,shift right=0.7ex,"{\pr_1^\ast}",swap] \ar[r,shift left=0.7ex,"\pr_0^\ast"]  &  \cG_1 \ar[r, shift left=1ex,"{\pr_{ij}^\ast}"] \ar[r,shift right=1ex] \ar[r] & \cG_2 
\end{tikzcd}\right]
\]
where $\cG_i = \prod_{s \in E_i} \cG_s$. Here there are three functors $v_i^\ast \colon \cG_0 \to \cG_2$, $i=0,1,2$. Each occurs in two ways as the composition $v_i^\ast \simeq  \pr_{jk}^\ast \circ \pr_\ell^\ast$ for all $j,k,\ell$ and $i = j$ for $\ell = 0$ while $i = k$ for $\ell = 1$. 
The $2$-complex $\cG_\bullet$ comes equipped with natural invertible transformations $v_i^\ast \simeq  \pr_{jk}^\ast \circ \pr_\ell^\ast$ as part of the datum implicitly. All these transformations will be denoted by $\alpha$ without causing too much confusion. 

\begin{ex}
\label{ex:FETdescent}
    Let $f\colon  X' \to X$ be a morphism of schemes. We set $X'' = X' \times_X X'$ and $X'''=X' \times_X X' \times_X X'$, and assume that $X'$, $X''$ and $X'''$ are locally connected, i.e.\ are the disjoint union of their connected components. Then 
    \[
    E_\bullet = \big[
    \begin{tikzcd} 
    \pi_0(X') & \ar[l,shift right=0.7ex] \ar[l,shift left=0.7ex] \pi_0(X'')  & \ar[l,shift right=1ex] \ar[l,shift left=1ex] \ar[l] \pi_0(X''')
    \end{tikzcd}
    \big]
    \]
    is a $2$-complex of sets. The fibred category $\FEt$ over schemes induces a fibred category over $E_\bullet$ that leads to the $2$-complex 
    \[
    \begin{tikzcd} 
    \FEt_{X'} \ar[r,shift left=0.7ex,"{\pr_0^\ast}"] \ar[r,shift right=0.7ex,"\pr_1^\ast",swap]   &  \FEt_{X''}  \ar[r, shift left=1ex,"{\pr_{ij}^\ast}"] \ar[r,shift right=1ex] \ar[r] & \FEt_{X'''} 
    \end{tikzcd}
    \]
    of multi-Galois categories. 
\end{ex}


%%%---------------------------------------
\subsubsection*{Descent} 
We now recall the definition of the category $\Desc(\cG_\bullet)$ of descent data as the equaliser category of a fibred category in Galois categories $\cG_\bullet \to E_\bullet$  (compare with \cite[Definition~4.1]{Stix2006:vanKampen}):
\begin{enumerate}
    \item 
    Objects of $\Desc(\cG_\bullet)$ are pairs $(X,\ph)$ with an object $X$ of $\cG_0$ and an isomorphism in $\cG_1$
    \[
    \ph \colon \pr_0^\ast(X) \isomto \pr_1^\ast(X).
    \]
    The isomorphism $\ph$ is referred to as the \textbf{gluing isomorphism}. 
    \item
    The isomorphism $\ph$ must satisfy the cocycle relation in $\cG_2$ (well defined if the necessary identifications are made using the natural isomorphisms $\alpha$):
    \[
    \pr_{02}^\ast(\ph) = \pr_{12}^\ast(\ph) \circ \pr_{01}^\ast(\ph).
    \]
    We refer to \cite[Definition~4.1]{Stix2006:vanKampen} for details.  
    \item 
    Morphisms $(X,\ph) \to (Y,\psi)$ in $\Desc(\cG_\bullet)$ are morphisms $X \to Y$ in $\cG_0$ that commute with $\ph$ and $\psi$ in $\cG_1$, i.e. we have
    \[
    \pr_1^\ast(f) \circ \ph = \psi \circ \pr_0^\ast(f).
    \]
\end{enumerate}

\begin{ex}
For the $2$-complex of Galois categories coming from the fibred category $\FEt$ over schemes as in \cref{ex:FETdescent} we introduce the short-hand notation $\Desc(X'/X, \FEt)$ for the category of descent data. 

Fix a scheme $S$ and denote the category of schemes over $S$ by $\Sch/S$. We can consider tame covers relative $S$ as a fibred category $\FEtt_{-/S} \to \Sch/S$ in view of the base change stability of tameness. 
For a map of $S$-schemes $X'\to X$ as in \cref{ex:FETdescent} we obtain a full subcategory of tame descent data
\[
    \Desc(X'/X, \FEtt_{-/S}) \subseteq \Desc(X'/X, \FEt).
\]
\end{ex}

\begin{prop}
\label{prop:DD is galois for connected E}
    Let $E_\bullet$ be a connected $2$-complex of sets and $\cG_\bullet \to E_\bullet$ a fibred category in Galois categories. Then $\Desc(\cG_\bullet)$ is a Galois category.
\end{prop}

\begin{proof}
    A proof can be found in 
    \cite[Propositions~2.4,~4.4]{Stix2006:vanKampen}. We sketch the crucial steps. Finite limits and colimits are constructed `simplexwise'.  
    An object $(X,\ph)$ consists of objects $X_v$ for every vertex $v \in E_0$ together with isomorphisms 
    \[
    \ph_e\colon  \pr_0^\ast(X_{\pr_0(e)}) \isomto \pr_1^\ast(X_{\pr_1(e)})
    \]
    in $\cG_e$ for every edge $e \in E_1$. Every object $X_v$ has a degree defined as the size of its fibre $F_v(X_v)$ under any fibre functor $F_v$ of $\cG_v$. The isomorphisms $\ph_e$ show that the degree is locally constant on $|E_\bullet|$, hence constant. Similarly, a map $f \colon (X,\ph) \to (Y, \psi)$ has fibrewise, i.e. for $F_v(f) \colon F_v(X_v) \to F_v(Y_v)$, image and preimage of locally constant size. In particular, $f$ is an isomorphism if and only if for one vertex $v$ the fibrewise map $F_v(f)$ is bijective. This shows that for any $v \in E_0$ the forgetful functor $\Desc(\cG_\bullet) \to \cG_v$ followed by a fibre functor of $\cG_v$ is a fibre functor for $\Desc(\cG_\bullet)$. The rest is formal. 
\end{proof}

\begin{defi}
    Let $\cC$ be a Galois category and $E_\bullet$ a connected $2$-complex of sets. We consider the constant fibred category $\cC \times E_\bullet \to E_\bullet$ with fibre $\cC$, i.e.~the pull-back functors $\partial^\ast\colon  \cC_s \to \cC_t$ are all the identity $\cC \to \cC$, and all natural transformations $\alpha$ are given by the identity.  
    
    Sending an object $Z$ of $\cC$ to the constant descent datum $(X,\ph)$ with $X_v = Z$ for all $v \in E_0$ and gluing isomorphism $\ph = \id$ defines a functor
    \[
    \cC \la \Desc(\cC \times E_\bullet).
    \]
    Since $E_\bullet$ is connected, this functor is fully faithful (but might fail to be an equivalence in general).

    A functor $f_0^\ast\colon  \cC \times E_\bullet \to \cG_\bullet$ over $E_\bullet$ induces a functor $\Desc(\cC \times E_\bullet) \to \Desc(\cG_\bullet)$; precomposing with the constant descent datum yields a functor
    \[
    f^\ast \colon \cC \la \Desc(\cG_\bullet).
    \]
    We say that $f_0^\ast$ satisfies \textbf{descent} (resp.\ \textbf{effective descent}) if $f^\ast$ is fully faithful (resp.\ an equivalence).

    In geometric settings (see the example below), when $f_0^*$ (and so $f^*$) come from some morphism $f$ of schemes or rigid spaces, we will also be saying that "$f$ is of (effective) descent".
\end{defi}

\begin{ex}
In our guiding \cref{ex:FETdescent} coming from a fibred category over schemes, the property of (effective) descent concerns properties of the functors
\begin{equation}
\label{eq:CechDescentSettingForFEt}
    f^\ast \colon \FEt_X \la \Desc(X'/X, \FEt) \quad \textrm{ and } \quad f^\ast \colon \FEtt_{X/S} \la \Desc(X'/X, \FEtt_{-/S})
\end{equation}
induced by a morphism of schemes $f\colon X'\to X$.
\end{ex}
We now gather some known results about descent concerning these categories.
\begin{fact}\label{fact:eff-decent-maps-for-FEt}
    Let $f \colon X' \to X$ be a surjective map of schemes.
\begin{enumerate}
    \item 
    (\cite[IX Corollaire 3.3+Proposition 2.4]{SGA1}) If $f$ is quasi-compact and universally submersive, then $f$ satisfies  descent for $\FEt$.
    \item 
    (\cite[Corollary~5.18]{RydhSubmersions}) If $f$ is  universally topologically submersive and $X, X'$ are noetherian, then $f$ satisfies effective descent for $\FEt$.
    \item 
    (\cite[IX Proposition~4.1]{SGA1}+\cite[IX Th\'eor\`eme~4.12]{SGA1}) When $f$ is an fpqc covering\footnote{Our presentation of effective descent is geared towards schemes with discrete $\pi_0$ -- this case is sufficient for our purposes. In general, fibre products of fpqc maps can often lead outside of this case, e.g.\ $\Spec(k^{\rm sep}\otimes_k k^{\rm sep})$.} or proper and of finite presentation, then $f$ satisfies effective descent for $\FEt$.
\end{enumerate}
\end{fact}
As  $\FEtt_{X/S} \subset \FEt_X$ and $\Desc(X'/X, \FEtt_{-/S}) \subset \Desc(X'/X, \FEt)$ are full subcategories, full faithfulness of functors is preserved and gives the following corollary.
\begin{lem}\label{lem:f-desc-for-FEt-and-FEtt}
    If $f \colon X' \to X$ satisfies descent for $\FEt$, then it satisfies descent for $\FEtt_{-/S}$. In particular, this holds in  the situations listed in \Cref{fact:eff-decent-maps-for-FEt}.
\end{lem}
Note that we do not claim \emph{effective} descent to be satisfied in the last corollary. In fact, the effective descent property for $\FEt^t_{-/S}$ will often fail even if it holds for $\FEt$.  


\begin{ex}
We consider the simplest example where the last assertion fails, maybe the embodiment of all examples: the Artin-Schreier cover $\wp \colon X' = \bA^1_K \to X = \bA^1_K$, $\wp(x) = x^p - x$ over $S = \Spec(K)$, a field of characteristic $p$. The Artin-Schreier cover considered as an object of $\FEt_X$ has a base change to $X'$ that splits completely. It thus yields a descent datum with respect to the fibred category $\FEtt_{-/S}$. This descent datum is not effective, for the only cover it can descend to is the Artin-Schreier cover itself, and that cover is not a tame cover.
\end{ex}

This problem can be fixed by adding extra assumptions that only hold in some specific situations.

\begin{lem}\label{lem:f-eff-desc-for-FEt-and-FEtt}
    Assume that $f\colon X'\to X$ satisfies effective descent for $\FEt$. Then it satisfies effective descent for $\FEt^t_{-/S}$ if and only if for every finite \'etale $Y\to X$ whose base change $Y'=Y\times_X X'\to X'$ is tame relative to $S$, we have that $Y\to X$ is tame relative to $S$. This holds when for every $x\in X$, there exists a point $x'\in X'$ such that $f(x')=x$ and $k(x')=k(x)$, or more generally such that $k(x')/k(x)$ is tame with respect to all valuations on $k(x')$ with centre on $S$.
\end{lem}

%%%---------------------------------------
\subsubsection*{Van Kampen formula} 
Let $\cG_\bullet \to E_\bullet$ be a fibred category in Galois categories over a $2$-complex of sets. For each simplex $s \in E_i$ we choose a base point (i.e.\ a fibre functor) $F_s \colon \cG_s \to \Sets$.
For $\partial(t) = s$ in $E_\bullet$ and the resulting exact functor $\iota_{st} \colon \cG_s \to \cG_t$ we choose a path, i.e. a natural isomorphism of fibre functors  $\gamma_{st} \colon F_t \circ \iota_{st} \to F_s$. This determines a group homomorphism 
\begin{equation}
\label{eq:boundary map 2-complex of pi1s}
    \partial_\ast \colon \pi_1(\cG_t ,F_t)\la \pi_1(\cG_s,F_s).
\end{equation}

Note that the datum of the natural transformations $\alpha$ that are implicit in $\cG_\bullet$ leads to a flag $vef$ in $E_\bullet$ consisting of a face $f \in E_2$, an edge $e = \pr_{jk}(f) \in E_1$ and a vertex $v = \pr_i(e) \in E_0$ (the indices $i,j,k$ are part of the datum!) to an element $\alpha_{vef} \in \pi_1(\cC_v,F_v)$ that makes composition commutative:
\[
\begin{tikzcd}
\pi_1(\cG_f,F_f) \ar[rr,"\pr_{jk,\ast}"] \ar[d,"v_{\ell,\ast}",swap] && \pi_1(\cG_e,F_e) \ar[d,"\pr_{i,\ast}"] \\
\pi_1(\cG_v,F_v) \ar[rr,"\alpha_{vef}(-)\alpha_{vef}^{-1}",swap] && \pi_1(\cG_v,F_v)
\end{tikzcd}
\]
where $\ell$ is the appropriate index that is determined by $v_\ell = \pr_{jk} \circ \pr_i$.
The above endows $E_\bullet$ with group data and allows one to define a Galois category of local systems, see \cite[\S4.3]{Stix2006:vanKampen}, equivalent to $\Desc(\cG_\bullet)$. The van Kampen formula 
\cite[Corollary~3.3]{Stix2006:vanKampen} computes the fundamental group of  $\Desc(\cG_\bullet)$ in terms of the $2$-complex $E_\bullet$ endowed with its group data. 

\begin{prop}
\label{prop:DD-fgfp}
    Let $\cG_\bullet \to E_\bullet$ be a category fibred in Galois categories over a connected $2$-complex. 
    \begin{enumerate}[ref=(\arabic*)]
        \item
        \label{propitem:DD-fg}
        Suppose that $\pi_1(\cG_v,F_v)$ is  finitely generated for all $v \in E_0$, and that $E_i$ is finite for $i \leq 1$. 
        Then $\pi_1(\Desc(\cG_\bullet),F_v)$ is finitely generated.
        \item
        \label{propitem:DD-fp}
        Suppose that $\pi_1(\cG_v,F_v)$ is finitely presented for all $v \in E_0$, that $\pi_1(\cG_e,F_e)$ is  finitely generated for all $e \in E_1$, and that $E_i$ is finite for all $i \leq 2$.
        Then $\pi_1(\Desc(\cG_\bullet),F_v)$ is finitely presented.
    \end{enumerate}
\end{prop}
\begin{proof}
This is stated in \cite[IX Corollaire~5.2,~5.3]{SGA1} with a proof that requires a ``common'' base point for all Galois categories involved and leaves the general case to the reader as an exercise. The remaining ``exercise'' can be found in the van Kampen formula in \cite[Corollary~3.3]{Stix2006:vanKampen}.  
\end{proof}

\begin{cor}
\label{cor:DD-fgfp}
       Let $\cG_\bullet \to E_\bullet$ be a fibred category over a connected $2$-complex, fibred in Galois categories. Let $\cC$ be a Galois category with fibre functor $F$, and let $f\colon  \cC \times E_\bullet \to \cG_\bullet$  be an exact functor over $E_\bullet$.  
       \begin{enumerate}[ref=(\arabic*)]
           \item 
            \label{coritem:DD-fg}
           If $f$ satisfies descent and $\cG_\bullet$ satisfies the assumptions of \cref{prop:DD-fgfp} \labelcref{propitem:DD-fg}, then $\pi_1(\cC,F)$ is finitely generated.
           \item 
            \label{coritem:DD-fp}
           If $f$ satisfies effective descent and $\cG_\bullet$ satisfies the assumptions of \cref{prop:DD-fgfp} \labelcref{propitem:DD-fp}, then $\pi_1(\cC,F)$ is finitely presented.
       \end{enumerate}
\end{cor}
\begin{proof}
    A choice of a path $F \simeq F_v \circ f$ induces a homomorphism $\pi_1(\Desc(\cG_\bullet,F_v)) \to \pi_1(\cC,F)$. In the first case, this homomorphism is surjective, and in the second case it is an isomorphism. Then apply \cref{prop:DD-fgfp}.
\end{proof}


%%%---------------------------------------
\subsubsection*{Graph of groups} 

In a $2$-complex of sets $E_\bullet$ the set $E_2$ might be empty. The $2$-complex in reality is then a graph. In this case the multi-Galois category $\cG_2$ for a fibred category in Galois categories is the category with just one object and one morphism. This leads in the definition of $\Desc(\cG_\bullet)$ to simplifications: the cocycle condition becomes automatic. 

We can truncate any given $2$-complex $E_\bullet$ by replacing $E_2$ with the empty set to obtain $E_{\le 1}$. This yields a right adjoint to the inclusion functor with counit $E_{\leq 1} \to E_\bullet$. Similarly, any fibred category $\cG_\bullet \to E_\bullet$ in Galois categories can be truncated by restriction to obtain a fibred $\cG_{\leq 1} \to E_{\leq 1}$, fibred in Galois categories. 

\begin{prop}
\label{prop:truncated DD}
    Let $\cG_\bullet \to E_\bullet$ be a fibred category in Galois categories over a connected $2$-complex.
    \begin{enumerate}
        \item 
        The natural forgetful functor `forget the cocycle condition'
        \[
        \Desc(\cG_\bullet) \la \Desc(\cG_{\leq 1})
        \]
        is fully faithful.
        \item
        A functor $\cC \times E_\bullet \to \cG_\bullet$ satisfies descent if and only if the restriction $\cC \times E_{\leq 1} \to \cG_{\leq 1}$ satisfies descent.
    \end{enumerate}
\end{prop}
\begin{proof}
    Obvious.
\end{proof}

Combining \cref{prop:truncated DD} with \cref{cor:DD-fgfp} shows that finite generation of fundamental groups depends only on truncation of descent data at level $1$. Furthermore, since an exact functor between Galois categories is fully faithful if and only if the induced homomorphism of fundamental groups is surjective \cite[V Proposition~6.9]{SGA1}, we can spell out concretely when an object $(X,\ph)$ of $\Desc(\cG_{\leq 1})$ is connected in terms of this equivalence. 

The underlying $2$-complex in the $1$-truncated setting is actually a graph $|E_{\leq 1}|$. It is endowed with profinite groups and thus becomes a graph of groups (cf.\ \cite[I.4.4 Definition 8]{Serre:trees}) as follows. For every simplex $s \in E_i$ we use the group $\pi_1(\cG_s,F_s)$ and for every edge $e \in E_1$ with boundary vertices $v_i = \pr_i(s) \in E_0$ we obtain group homomorphisms
\[
\pi_1(\cG_{v_0},F_{v_0}) \xleftarrow{\pr_{0,\ast}} 
\pi_1(\cG_e,F_e)
\xrightarrow{\pr_{1,\ast}} \pi_1(\cG_{v_1},F_{v_1}) 
\]
where we have implicitly also fixed paths for change of base points as in \labelcref{eq:boundary map 2-complex of pi1s}. A \textbf{local system} $M = (M_s)_{s \in E_{\leq 1}}$ on this graph of groups now consists of a finite set $M_s$ for every vertex $s \in E_i$, $i=0,1$ with a continuous action by $\pi_1(\cG_s,F_s)$ together with, for every map $\partial \colon e \to v$ in the category $E_{\leq 1}$ a monodromy map $m_{\partial} \colon M_e \to M_v$ which is bijective and equivariant along the homomorphism $\partial_\ast \colon \pi_1(\cG_e,F_e) \to \pi_1(\cG_v,F_v)$. Local systems form a category with the obvious notion of morphism. If, moreover, $|E_{\leq 1}|$ is connected, then local systems form a Galois category equivalent to $\Desc(\cG_{\leq 1})$ by \cite[Proposition~4.4]{Stix2006:vanKampen}. 

Let $T$ be a maximal tree in the graph $|E_{\leq 1}|$. The functor $F_T$ on $\Desc(\cG_{\leq 1})$ that assigns the set of connected components of the restriction to $T$ of the underlying local system $M = (M_s)_{s \in E_{\leq 1}}$ yields a group homomorphism $\pi_1(|E_{\leq 1}|,T)^\wedge \to \pi_1(\Desc(\cG_{\leq 1}),F_T)$. Here $\pi_1(|E_{\leq 1}|,T)^\wedge$ denotes the profinite completion of the topological fundamental group of the graph $|E_{\leq 1}|$ with the maximal tree $T$ as a base point. Furthermore, for every vertex $v \in E_0$, the inclusion $v \in T$ induces a natural transformation of fibre functors $F_v \isomto F_T$. This leads to group homomorphisms $\pi_1(\cG_v,F_v) \to \pi_1(\Desc(\cG_{\leq 1}),F_T)$ that together combine to a natural group homomorphism
\[
    \bigg(\underset{v \in E_0}{\resizebox{.3cm}{!}{$\ast$}} \pi_1(\cG_v,F_v)\bigg)  \ast \pi_1(|E_{\leq 1}|,T)^\wedge \surj \pi_1(\Desc(\cG_{\leq 1}))  .
\]
By \cite[Proposition~2.7]{Stix2006:vanKampen} this is surjective. In particular, this proves \cref{prop:DD-fgfp} \labelcref{propitem:DD-fg}.

To a local system $M$ we can associate a graph $M_{\leq 1} = \big[\begin{tikzcd}
M_0 & \ar[l,shift right=0.7ex] \ar[l,shift left=0.7ex]  M_1 
\end{tikzcd}\big]$ over $E_{\leq 1}$ by 
\[
M_{\leq 1} = 
\Big[\begin{tikzcd}
\amalg_{v \in E_0} M_v & \ar[l,shift right=0.7ex] \ar[l,shift left=0.7ex] \amalg_{e \in E_1} M_e 
\end{tikzcd}\Big] 
\la
E_{\leq 1} = \Big[\begin{tikzcd}
E_0 & \ar[l,shift right=0.7ex] \ar[l,shift left=0.7ex] E_1 
\end{tikzcd}\Big]
\]
with source map $\pr_0$ and target  map  $\pr_1$ defined on $M_e$ by the monodromy maps
\[
M_{\pr_0(e)} \xleftarrow{m_{\pr_0}} M_e  \xrightarrow{m_{\pr_1}} M_{\pr_1(e)}
\]
of the local system $M$.
Recall that the $M_s$ for $s \in E_i$ are $\pi_1(\cG_s,F_s)$-sets and that the set of connected components $\pi_0(M_s)$ is the set of orbits for this group action. Associated to the graph $M_{\leq 1}$ there is the graph 
\[
\pi_0(M_{\leq 1}) = 
\Big[\begin{tikzcd}
\amalg_{v \in E_0} \pi_0(M_v) & \ar[l,shift right=0.7ex] \ar[l,shift left=0.7ex] \amalg_{e \in E_1} \pi_0(M_e) 
\end{tikzcd}\Big]
\]
of connected components. We have a factorisation
\[
M_{\leq 1} \la \pi_0(M_{\leq 1}) \la E_{\leq 1}.
\]

\begin{lem} \label{lem:local-system-connected}
    The object of $\Desc(\cG_{\leq 1})$ corresponding to the local system $M$ on the associated graph of groups is connected if and only if the associated graph $\pi_0(M_{\leq 1})$ of connected components is connected.
\end{lem}

\begin{proof}
    Since $\Desc(\cG_{\leq 1})$ is equivalent to the Galois category of local systems, we may focus on connectedness of $M$. 
    Passing from $M$ to $\pi_0(M_{\leq 1})$ is a functor that preserves disjoint unions, and $\pi_0(M_{\leq 1})$ is empty if and only if $M$ is empty. Conversely, any decomposition $M = A \sqcup B$ arises by a decomposition $M_s = A_s \sqcup B_s$ as $\pi_1(\cG_s,F_s)$-sets for every simplex $s$ of $E_{\leq 1}$ that are compatible with monodromy operators.  
    This always arises as the preimage of a decomposition $\pi_0(M_s) = \ov{A}_s \sqcup \ov{B}_s$ for every simplex $s$ of $E_{\leq 1}$. 
    But this agrees with the datum of a decomposition of 
    $\pi_0(M_{\leq 1})$ as a graph.
\end{proof}


%%%---------------------------------------
\subsection{Finite generation and finite presentation}
%%%---------------------------------------

We end \cref{sec:pit schemes} with a general finite generation result for the tame fundamental group of connected varieties over an algebraically closed field. The proof proceeds by reduction to the case of a variety admitting a smooth projective snc compactification, in which case the result is known due to \cite[Theorem~1.1]{EsnaultKindler2016:LefschetzTheoremsTamely} and the case of curves \cite[XIII, Corollaire 2.12]{SGA1}.
We also give a sufficient condition for finite presentation which relies on \cite[Theorem~1.1]{EsnaultShustermanSrinivas2022:FinitePresentationTame} instead. 

\begin{thm} \label{thm:naked pi1tame fg}
    Let $X$ be a connected scheme of finite type over an algebraically closed field $k$. Then $\pitame(X/k,\bar x)$ is finitely generated. 
\end{thm}

\begin{proof}
Let $h_1 \colon X_1 \to X$ be the map from the disjoint union $X_1 = \coprod U$ of an affine open covering of $X$ by finitely many affine Zariski open $U$. 
Let further $h_2\colon X_2 \to X_1$ be the normalisation map of $X_1$, and let $X_2 = \coprod Z_\alpha$ be the decomposition into (finitely many) connected components. 
By \cite[Theorem~4.1]{deJong1996:SmoothnessSemistabilityAlterations} there exist  alterations $f_\alpha\colon  Y_\alpha \to Z_\alpha$ with a smooth projective compactification
$j\colon  Y_\alpha \inj \ov{Y}_\alpha$ such that the boundary $D_\alpha = \ov{Y}_\alpha - Y_\alpha$ is normal crossing. 
Set $Y$ equal to $\coprod_\alpha Y_\alpha$ and let $h_3: Y \to X_2$ be the map $\amalg f_\alpha$. 
The composition $h = h_1\circ  h_2 \circ h_3 \colon Y \to X$ is quasi-compact and universally submersive. So $h$ is of effective descent for $\FEt$ by \cref{fact:eff-decent-maps-for-FEt}.
As observed in \Cref{lem:f-desc-for-FEt-and-FEtt}, it follows that $h$ is of  descent for $\FEtt_{-/k}$.
Therefore the map
\[
h_\ast\colon  \pi_1\big(\Desc(X^\nu/X,\FEtt_{-/k}), F_{\bar y}\big) \surj \pitame(X/k,\bar x)
\]
is surjective (where $F_{\bar y}$ is the fibre functor at a geometric point $\bar y$ of $X^\nu$ above $\bar x$). 
By \cite[Proposition~3.2]{EsnaultShustermanSrinivas2022:FinitePresentationTame}, the group $\pitame(Y_\alpha/k,\bar y) = \pitame(\ov{Y}_\alpha,D_\alpha; \bar y)$ is finitely generated for every $\alpha$.  
As $\pi_0(Y)$ and $\pi_0(Y \times_X Y)$ are finite sets,
it follows by \cref{cor:DD-fgfp} \labelcref{coritem:DD-fg}
that $\pitame(X/k,\bar x)$ is finitely generated. 
\end{proof}

\cref{thm:naked pi1tame fg} recovers the finite generation part of \cite[Th\'eor\`eme~4.11]{Orgogozo2003:AlterationsGroupeFondamental}:

\begin{cor}
    The prime-to-$p$ quotient of the \'etale/tame fundamental group $\pi_1(X)^{\wedge p'}$ of any connected scheme of finite type over an algebraically closed field of char $p>0$ is finitely generated.
\end{cor}

\begin{proof}
    The category of all étale covers of~$X$ whose connected components have Galois closures of degrees prime to~$p$ is a full subcategory of $\FEtt_{X/K}$ satisfying the conditions of \cref{lem:Galois-subcategory-criterion}.
    The corresponding fundamental group is thus a quotient of $\pitame(X/K)$ and inherits the property of being finitely generated.
\end{proof}

\begin{thm} \label{thm:naked pi1tame fp}
    Let $X$ be a connected scheme of finite type over an algebraically closed field $k$. Suppose that there exists a morphism of finite type $f\colon Y\to X$ such that 
    \begin{enumerate}[(i)]
        \item 
        $f$ is universally topologically submersive (e.g.\ surjective and either proper or flat);
        
        \item 
        for every $x\in X$ there exists a point $y\in Y$ such that $f(y)=x$ and $k(y)=k(x)$, or more generally such that the extension $k(y)/k(x)$ is tame with respect to all valuations on $k(x)$ which are trivial on $k$;
        
        \item 
        there exists a projective snc compactification $Y \hookrightarrow \overline{Y}$ over $k$.
    \end{enumerate}
     Then $\pitame(X/k)$ is finitely presented.
\end{thm}

\begin{proof}
Condition (i) and \Cref{fact:eff-decent-maps-for-FEt} guarantee that the morphism $f$ satisfies effective descent for finite \'etale covers. Condition (ii) and \Cref{lem:f-eff-desc-for-FEt-and-FEtt} ensure that the same is true for $\FEtt_{-/k}$, i.e.\ finite \'etale covers tame over $k$. For every connected component $Z$ of $Y$, $Y\times_X Y$, or $Y\times_X Y\times_X Y$, the group $\pitame(Z/k)$ is finitely generated by \cref{thm:naked pi1tame fg}. Moreover, condition (iii) and \cite[Theorem~1.1]{EsnaultShustermanSrinivas2022:FinitePresentationTame} together imply that for every connected component $Z$ of $Y$, the group $\pitame(Z/k)$ is finitely presented. We are now able to apply \cref{cor:DD-fgfp}~\labelcref{coritem:DD-fp}.
\end{proof}

\begin{rmk}
    We mention a few further results for variants of tame fundamental groups of schemes.
    \begin{itemize}
        \item 
        The K\"unneth formula holds for $\pitame$ if both factors are smooth connected varieties over an algebraically closed field and $\pitame$ is computed with respect to curve tame covers. This was recently proved by de~Jong in \cite[Theorem A.1]{deJong2026:KuennethForPitameAppendixJLR2026} and was previously known for varieties with good compactifications (the boundary is a normal crossing divisor) by  
        \cite[Th\'eor\`eme~5.1]{Orgogozo2003:AlterationsGroupeFondamental}. 
        \item 
        Invariance of $\pitame$ under geometric (separably closed) base change of base fields for varieties with good compactification: \cite[Remarque~5.3]{Orgogozo2003:AlterationsGroupeFondamental}. See also \cite{Landesman}.
    \end{itemize}
\end{rmk}
